% !PRINTED 677
\weberpage{677}
This now proves the theorem:
% [source note] The source prints der Satz; den Satz is apparently intended, since the preceding argument establishes the theorem which follows. The English renders that intended direction; the German preserves the printed wording.

4. \gesperrt{If $g_1, g_2, \ldots, g_\nu$ is an arbitrary system of real numbers that are not all equal to one another, then in every algebraic field $\Omega$ a unit $\varepsilon$ can be so determined that the sum}
\[
g_1 l_1 + g_2 l_2 + \cdots + g_\nu l_\nu
\]
\gesperrt{is different from zero, when $l_s$ is the system of the conjugate logarithms of $\varepsilon$.}

This important theorem, which is the foundation for what follows, goes back, though in a somewhat altered version, to \gesperrt{Dirichlet}. This formulation was given by \gesperrt{Minkowski}.

\begin{center}
\S.\,186.\\
\gesperrt{Systems of independent units and systems of exponents of the units.}
\end{center}

If in the theorem last proved we set $g_1 = 1, g_2 = 0, \ldots, g_\nu = 0$, which is permissible once we assume from now on that $\nu > 1$, it follows that there is in $\Omega$ a unit $\varepsilon$ for which one of the conjugate logarithms, $l_1$, is different from zero. Since every power of $\varepsilon$ has the same property, there are also infinitely many such units.

This theorem now admits a very important generalisation:

5. \gesperrt{If $s \leqq \nu - 1$, then in the field $\Omega$ a system of units with the corresponding conjugate logarithms}
\weq{1}{
\begin{array}{l}
\varepsilon_1 : l_{1,1}, l_{1,2}, \ldots, l_{1,\nu} \\
\varepsilon_2 : l_{2,1}, l_{2,2}, \ldots, l_{2,\nu} \\
\cdots\cdots\cdots\cdots\cdots\cdots \\
\varepsilon_s : l_{s,1}, l_{s,2}, \ldots, l_{s,\nu}
\end{array}
}
\gesperrt{can be so determined that an arbitrary one of the $s$-rowed determinants of the matrix of the $l_{h,k}$, say}
\[
L_s = \Sigma \pm l_{1,1} l_{2,2} \ldots l_{s,s}
\]
\gesperrt{is different from zero.}

For $s = 1$ the theorem as stated is correct by the remark made at the outset. We therefore assume that it
%CONT>>
